\section[Hex - Dez]{Hex $\to$ Dez}
\begin{tabular}{l|cccccc}
    \textbf{Hex} & A & B & C & D & E & F \\
    \textbf{Dec} & 10 & 11 & 12 & 13 & 14 & 15 \\
    \textbf{Bin} & 1010 & 1011 & 1100 & 1101 & 1110 & 1111
\end{tabular}

\section{Instruction Set Architecture (ISA)}
\label{subsec:isa_differences}

L'evoluzione delle architetture ARM ha portato alla coesistenza di diversi set di istruzioni.

\begin{itemize}
    \item \textbf{ARM (32-bit) (Obsoleto):} Istruzioni a lunghezza fissa di 32 bit. 
    Garantiscono massime prestazioni ma occupano maggior spazio in memoria.
    \item \textbf{Thumb (16-bit):} Sottoinsieme del set ARM con istruzioni a 16 bit. 
    Ideale per ridurre il footprint della memoria (circa 30\% di risparmio), 
    a scapito di una minore flessibilità operativa.
    \item \textbf{Thumb-2 (16/32-bit):} Tecnologia a lunghezza variabile che combina 
    istruzioni a 16 e 32 bit. Permette di ottenere una densità di codice simile a Thumb 
    con prestazioni analoghe al set ARM a 32 bit.
\end{itemize}

\subsubsection{Istruzioni Comuni}
\begin{tabularx}{\linewidth}{|l|X|}
    \hline
    \textbf{Mov/Load} & \texttt{MOV} (trasferimento), \texttt{LDR/STR} (RAM/ROM access), \texttt{PUSH/POP} (Stack). \\ \hline
    \textbf{Bitwise} & \texttt{ANDS} (\&), \texttt{ORRS} (|), \texttt{EORS} (XOR), \texttt{BICS} (Bit Clear). \\ \hline
    \textbf{Arithmetic} & \texttt{ADDS}, \texttt{SUBS}, \texttt{MULS}. \\ \hline
    \textbf{Shifts} & \texttt{LSLS} (Logical Left), \texttt{LSRS} (Logical Right), \texttt{ASRS} (Arithmetic Right), \texttt{RORS} (Rotate). \\ \hline
\end{tabularx}

\section{Logica della ALU e Program Status Register}
\label{sec:alu_logic}

L'Unità Aritmetico-Logica (ALU) opera esclusivamente su stringhe binarie di bit, 
ignorando la semantica del segno (interpretazione \textit{signed} vs \textit{unsigned}).
Tale agnosticismo è reso possibile dalla rappresentazione in complemento a due, che 
permette di unificare le operazioni di addizione e sottrazione.

\subsection{Meccanismo delle Flag (APSR)}
\label{subsec:apsr_flags}

Al termine di ogni operazione aritmetica, la ALU aggiorna meccanicamente quattro flag 
fondamentali nel \textit{Application Program Status Register} (APSR). 
Queste flag descrivono le proprietà del risultato binario grezzo:

\begin{itemize}
    \item \textbf{N (Negative):} Riflette il bit più significativo (MSB) del risultato. Se $MSB=1$, la flag è impostata a 1.
    \item \textbf{Z (Zero):} Impostata a 1 se tutti i bit del risultato sono pari a zero.
    \item \textbf{C (Carry):} Indica un riporto oltre il bit più significativo. 
\begin{itemize}
    \item Fondamentale per operazioni \textit{unsigned}. \crd{Invertito per la sottrazione!}
    \item \textbf{Meccanismo \crd{sottrazione}:} $(A + \text{NOT}(B) + 1)$. \quad Zweierkompl.
    \begin{itemize}
        \item $C = 1 \implies A \ge B$ (Nessun prestito / Not-Borrow).
        \item $C = 0 \implies A < B$ (Prestito avvenuto / Borrow).
    \end{itemize}
\end{itemize}
    \item \textbf{V (oVerflow):} Indica un errore di segno nel calcolo in complemento 
    a due, Fondamentale per operazioni \textit{signed}.
    \begin{itemize}
        \item $(+) + (+) = (-)$ Impossibile
        \item $(-) + (-) = (+)$ Impossibile
    \end{itemize}
\end{itemize}

\subsubsection{Interpretazione del Segno e Salti Condizionati}
\label{subsec:sign_interpretation}

Poiché la ALU aggiorna simultaneamente sia la flag $C$ che la flag $V$,
la distinzione tra variabile \texttt{int} e \texttt{uint} non risiede nell'hardware, 
ma nel software.

